% ex: ts=2 sw=2 sts=2 et filetype=tex
% SPDX-License-Identifier: CC-BY-SA-4.0

\documentclass[10pt,addpoints]{exam}

\usepackage[utf8]{inputenc}
\usepackage[T1]{fontenc}
\usepackage[letterpaper,top=50mm]{geometry}
\usepackage{graphicx}
\usepackage[table]{xcolor}
%\usepackage[spanish]{babel}
\usepackage{multicol}
\usepackage{amsmath}
\usepackage{listings}

\usepackage{tikz}
\newcommand{\pythagwidth}{3cm}
\newcommand{\pythagheight}{2cm}

\pagestyle{headandfoot}
%\headrule
\lhead{\includegraphics[scale=0.10]{img/tsj.png}}
\rhead{\textbf{Diseño y Análisis de Algoritmos}}
\lfoot{Nombre:~\makebox[4.5in]{\hrulefill} }
\cfoot{}
\rfoot{Página \thepage\ de \numpages}

\pointpoints{punto}{puntos}
\renewcommand{\solutiontitle}{\textbf{Solución: }}
\renewcommand{\theenumi}{\alph{enumi})}

%\printanswers

\begin{document}
\begin{center}
  \sffamily\textbf{EXAMEN II}
\end{center}
\begin{flushright}
9 de abril de 2024
\end{flushright}

\begin{questions}
\begin{EnvFullwidth}
  \sffamily\textbf{Sección I.- Completar los siguientes enunciados.
  Valor: 60 puntos (5 punto c/u).}
\end{EnvFullwidth}

% ex: ts=2 sw=2 sts=2 et filetype=tex
% SPDX-License-Identifier: CC-BY-SA-4.0

\question Función que devuelve una secuencia de números:
          \fillin[range()]

% ex: ts=2 sw=2 sts=2 et filetype=tex
% SPDX-License-Identifier: CC-BY-SA-4.0

\question Manejo de espacios al principio de la linea para delimitar
          bloques de código:
          \fillin[Sangría]


% ex: ts=2 sw=2 sts=2 et filetype=tex
% SPDX-License-Identifier: CC-BY-SA-4.0

\question Función que divide una cadena en una lista especificando un
          separador:
          \fillin[split()]

% ex: ts=2 sw=2 sts=2 et filetype=tex
% SPDX-License-Identifier: CC-BY-SA-4.0

\question \textbf{join()}: Función que retorna una cadena formada por los
          elementos de una lista separados por un delimitador.

% ex: ts=2 sw=2 sts=2 et filetype=tex
% SPDX-License-Identifier: CC-BY-SA-4.0

\question  

% ex: ts=2 sw=2 sts=2 et filetype=tex
% SPDX-License-Identifier: CC-BY-SA-4.0

\question Se constituye por un identificador, un tipo de dato, un valor y un
          ámbito:
          \fillin[Variable]

% ex: ts=2 sw=2 sts=2 et filetype=tex
% SPDX-License-Identifier: CC-BY-SA-4.0

\question \tf[V] El modo "r" abre un archivo de solo lectura y envía mensaje
          de error si el archivo no existe.


% ex: ts=2 sw=2 sts=2 et filetype=tex
% SPDX-License-Identifier: CC-BY-SA-4.0

\question Función que divide una cadena en una lista especificando un
          separador:
          \fillin[split()]

% ex: ts=2 sw=2 sts=2 et filetype=tex
% SPDX-License-Identifier: CC-BY-SA-4.0

\question \textbf{join()}: Función que retorna una cadena formada por los
          elementos de una lista separados por un delimitador.

% ex: ts=2 sw=2 sts=2 et filetype=tex
% SPDX-License-Identifier: CC-BY-SA-4.0

\question  

% ex: ts=2 sw=2 sts=2 et filetype=tex
% SPDX-License-Identifier: CC-BY-SA-4.0

\question Se constituye por un identificador, un tipo de dato, un valor y un
          ámbito:
          \fillin[Variable]

% ex: ts=2 sw=2 sts=2 et filetype=tex
% SPDX-License-Identifier: CC-BY-SA-4.0

\question \tf[V] El modo "r" abre un archivo de solo lectura y envía mensaje
          de error si el archivo no existe.

\end{questions}

\begin{questions}
\begin{EnvFullwidth}
  \sffamily\textbf{Sección II.- Resolver el siguiente problemas utilizando
  instrucciones del lenguaje Python. Valor 40 puntos.}
\end{EnvFullwidth}

% ex: ts=2 sw=2 sts=2 et filetype=tex
% SPDX-License-Identifier: CC-BY-SA-4.0

\question Función que divide una cadena en una lista especificando un
          separador:
          \fillin[split()]


\end{questions}

\end{document}
